% ============================================================
%  SEÇÃO 4 — CRONOGRAMA
%
%  O cronograma mostra o planejamento do trabalho de PS I a PS II.
%  Deve cobrir todas as etapas até a defesa final.
%
%  --> Ajuste os meses e as atividades conforme seu projeto.
%     Marque com X os meses em que cada atividade ocorrerá.
% ============================================================

\chapter{CRONOGRAMA}
\label{cap:cronograma}

O cronograma apresentado no Quadro~\ref{qua:cronograma} organiza as etapas de execução deste trabalho ao longo dos semestres de 2026.1 (Projeto Supervisionado I) e 2026.2 (Projeto Supervisionado II).

\begin{table}[htb]
\caption{Cronograma de execução do trabalho}
\label{qua:cronograma}

% ----------------------------------------------------------
%  LEGENDA: $\checkmark$ = atividade prevista para o mês
%  Ajuste as colunas de mês conforme seu semestre
% ----------------------------------------------------------

\renewcommand{\arraystretch}{1.4}
\begin{tabularx}{\textwidth}{|X|c|c|c|c|c|c|c|c|c|c|}
  \hline
  \multirow{2}{*}{\textbf{Atividade}} &
  \multicolumn{5}{c|}{\textbf{2026.1 (PS I)}} &
  \multicolumn{5}{c|}{\textbf{2026.2 (PS II)}} \\
  \cline{2-11}
  & \textbf{Fev} & \textbf{Mar} & \textbf{Abr} & \textbf{Mai} & \textbf{Jun}
  & \textbf{Ago} & \textbf{Set} & \textbf{Out} & \textbf{Nov} & \textbf{Dez} \\
  \hline

  % ----------------------------------------------------------
  % %->  ATIVIDADES DE PS I — preencha com $\checkmark$ nos meses corretos
  % ----------------------------------------------------------

  Definição do tema e problema de pesquisa
  & $\checkmark$ & $\checkmark$ & & & & & & & & \\
  \hline

  Revisão bibliográfica e referencial teórico
  & & $\checkmark$ & $\checkmark$ & $\checkmark$ & & & & & & \\
  \hline

  Redação da metodologia
  & & $\checkmark$ & $\checkmark$ & & & & & & & \\
  \hline

  Escrita e formatação do Projeto de Pesquisa
  & & & $\checkmark$ & $\checkmark$ & & & & & & \\
  \hline

  Entrega e defesa do Projeto de Pesquisa (PS I)
  & & & & & $\checkmark$ & & & & & \\
  \hline

  % ----------------------------------------------------------
  % %->  ATIVIDADES DE PS II — ajuste conforme seu tipo de projeto
  % ----------------------------------------------------------

  % Para projeto de SOFTWARE — adapte conforme necessário:
  Levantamento de requisitos
  & & & & & & $\checkmark$ & & & & \\
  \hline

  Modelagem e prototipagem (Figma / UML)
  & & & & & & $\checkmark$ & $\checkmark$ & & & \\
  \hline

  Implementação do sistema
  & & & & & & & $\checkmark$ & $\checkmark$ & & \\
  \hline

  Testes e validação
  & & & & & & & & $\checkmark$ & $\checkmark$ & \\
  \hline

  % Para projeto de PESQUISA — substitua as linhas acima por:
  % Coleta de dados
  % & & & & & & $\checkmark$ & $\checkmark$ & & & \\
  % \hline
  % Análise dos dados
  % & & & & & & & $\checkmark$ & $\checkmark$ & & \\
  % \hline

  Redação do TCC completo
  & & & & & & & & $\checkmark$ & $\checkmark$ & \\
  \hline

  Revisão e formatação final (ABNT)
  & & & & & & & & & $\checkmark$ & \\
  \hline

  Entrega e defesa do TCC (PS II)
  & & & & & & & & & & $\checkmark$ \\
  \hline

\end{tabularx}
\fonte{Elaborado pelo autor (\oano).}
\end{table}

% ----------------------------------------------------------
% %->  DESCRIÇÃO DAS ETAPAS (opcional, mas recomendado)
%     Explique brevemente cada etapa do cronograma.
% ----------------------------------------------------------

As etapas previstas no cronograma são descritas a seguir:

\begin{description}

  \item[\textbf{Definição do tema e problema de pesquisa}]
    Inclui a escolha do tema, formulação do problema de pesquisa,
    definição dos objetivos e justificativa, com validação pelo orientador.

  \item[\textbf{Revisão bibliográfica e referencial teórico}]
    Levantamento sistemático de artigos e obras nas bases IEEE Xplore,
    ACM Digital Library, Periódicos CAPES e Google Scholar,
    organizados com Zotero.

  \item[\textbf{Redação da metodologia}]
    Definição e documentação do tipo de pesquisa, abordagem,
    procedimentos e instrumentos de coleta de dados.

  \item[\textbf{Escrita e formatação do Projeto de Pesquisa}]
    Redação e formatação do documento completo em \LaTeX\ (Overleaf)
    seguindo as normas ABNT (NBR 14724:2011).

  \item[\textbf{Entrega e defesa do Projeto de Pesquisa (PS I)}]
    Apresentação oral de 12--15 minutos perante banca examinadora,
    conforme Art. 36º do Regulamento Institucional de TCC.

  % Adicione descrições das etapas de PS II conforme seu projeto

\end{description}
